\documentclass[12pt]{article}
\usepackage[utf8]{inputenc}
\usepackage[T1]{fontenc}
\usepackage{hyperref}
\title{Feroge Language Resources for K-12 Curriculum Development}
\author{Compiled by Claude Sonnet 5.5 (claude-sonnet-5-5) for Andr\'as Kornai}
\date{}
\begin{document}
\maketitle
\section{Introduction}
Feroge (also Kaligi, Kalige; ISO 639-3 fer; Glottolog fero1244) is spoken in Western Bahr el Ghazal, South Sudan. Search results classify it as Ubangian (Adamawa-Ubangi). Speaker estimates conflict: about 26,000 (Wikipedia, citing Ethnologue) versus 8,000 (Global Recordings Network). The language is reported as unwritten. Very little exists.
\section{Primary Educational Resources}
None found. No primers, textbooks or school programmes appeared in the search results.
\section{Linguistic and Reference Materials}
\begin{itemize}
\item Wikipedia, Feroge/Kaligi: \url{https://en.wikipedia.org/wiki/Feroge}
\item Wikipedia, Kaligi language: \url{https://en.wikipedia.org/wiki/Kaligi_language}
\item Numerals 1--10 in Feroge (Max Planck numerals database; the contributor notes that further data is needed): \url{https://lingweb.eva.mpg.de/channumerals/Feroge.htm}
\item Wikitongues language page (listed as unwritten, EGIDS 7): \url{https://wikitongues.org/languages/fer/}
\end{itemize}
\section{Religious and Community Resources}
\begin{itemize}
\item Joshua Project profile of Kaligi/Feroge: \url{https://joshuaproject.net/people_groups/12451}
\end{itemize}
\section{Digital and Audio Resources}
Global Recordings Network reports no recordings available: \url{https://globalrecordings.net/en/language/fer}
\section{Curriculum Development Context}
There is no known writing system, orthography or school use. Any curriculum would first require orthography development and documentation with the community.
\end{document}